\section{Experiments}
\label{sec:experiments}

To evaluate whether phased orchestration improves task completion quality, we design four experiments that collectively address the concerns raised in prior review: (1) an ablation study isolating each phase's contribution, (2) a control marker effectiveness study, (3) a three-protocol consistency evaluation, and (4) a cost-benefit analysis. All experiments are conducted through the Pigs HTTP API, requiring only API access---no GPU resources.

\subsection{Experimental Setup}

\textbf{System configuration.} All experiments use the Pigs API server with default orchestration limits (\texttt{max\_post\_iterations=3}, \texttt{max\_pre\_replans=2}), continuation TTL of 30 minutes, capacity of 256, and Chinese (zh) phase prompts. Thinking effort is set to xhigh for OpenAI and max for Anthropic.

\textbf{Benchmarks.} We use three public benchmarks:

\begin{center}
\begin{tabular}{lllr}
\toprule
Benchmark & Domain & Tasks & Public \\
\midrule
SWE-bench Lite~\citep{jimenez2024swebench} & Software engineering & 300 & \checkmark \\
GAIA~\citep{mialon2023gaia} & General AI assistant & 300 & \checkmark \\
AgentBoard~\citep{ma2024agentboard} & Multi-turn agent & 9 tasks & \checkmark \\
\bottomrule
\end{tabular}
\end{center}

\textbf{Ablation conditions.} For Experiment 1, we define four conditions:
\begin{itemize}[nosep]
    \item \textbf{Condition A (Baseline)}: Direct API call without orchestration (non-\texttt{-pig} model, passthrough).
    \item \textbf{Condition B (Full)}: Complete Pre$\rightarrow$Executor$\rightarrow$Post orchestration (\texttt{-pig} model).
    \item \textbf{Condition C (Pre-only)}: Only the Pre phase (planning without execution verification).
    \item \textbf{Condition D (Executor+Post)}: Skip Pre, directly execute and review.
\end{itemize}

\textbf{Metrics.} We measure: (1) task completion rate (benchmark-specific), (2) average LLM API calls per task, (3) total input + output tokens per task, (4) wall-clock latency, and (5) error recovery rate (percentage of tasks where PIGFAIL triggered a successful replan).

\subsection{Experiment 1: Ablation Study}

\textbf{Setup.} We evaluate four conditions (A: baseline, B: full orchestration, C: Pre-only, D: Executor+Post) on SWE-bench Lite (300 tasks) and GAIA (300 tasks). Each task is sent as a \texttt{-pig} model request (Condition B/C/D) or a passthrough request (Condition A) to the Pigs API.

\textbf{Procedure.} For each task:
\begin{enumerate}[nosep]
    \item Send the task prompt to the Pigs API endpoint.
    \item If tool calls are returned, execute them locally and send results back.
    \item Repeat until the API returns a final response (no tool calls).
    \item Evaluate the final response against the benchmark's ground truth.
\end{enumerate}

\textbf{Hypothesis.} Full orchestration (B) outperforms baseline (A) on task completion rate, with the Pre phase contributing planning and the Post phase contributing verification.

\subsection{Experiment 2: Control Marker Effectiveness}

\textbf{Setup.} We compare two configurations on GAIA: (A) with \texttt{PIGFAIL} enabled (failure triggers replan), (B) with \texttt{PIGFAIL} disabled (failure treated as completion).

\textbf{Procedure.} We count: (1) tasks where PIGFAIL was emitted, (2) tasks where replan led to successful completion, (3) tasks that failed despite replan.

\textbf{Hypothesis.} PIGFAIL-based replan improves outcomes on tasks where the initial approach was flawed.

\subsection{Experiment 3: Three-Protocol Consistency}

\textbf{Setup.} We send the same set of 100 tasks through three API protocols: OpenAI Chat (\texttt{/chat/completions}), Anthropic Messages (\texttt{/v1/messages}), and OpenAI Responses (\texttt{/responses}).

\textbf{Procedure.} For each task and protocol, we record: (1) whether the task completed, (2) the visible text output, (3) any tool calls made. We measure consistency as the percentage of tasks where all three protocols agree on completion status.

\textbf{Hypothesis.} The protocol-native request preservation ensures consistent behavior across all three protocols.

\subsection{Experiment 4: Cost-Benefit Analysis}

\textbf{Setup.} On SWE-bench Lite, we measure per-task: (1) total input tokens, (2) total output tokens, (3) number of API calls, (4) wall-clock latency, (5) task completion.

\textbf{Procedure.} We compare Condition A (baseline) vs.\ Condition B (full orchestration) and compute the overhead ratio for each metric.

\textbf{Hypothesis.} The token overhead of orchestration (3 phases vs.\ 1) is justified by improved completion rate.
